% !TeX root = ../../Book3.tex
\section{单连通区域}
\label{b3:sec:0501}

把一个区域映成圆盘，是寻找全局单射函数的问题。局部逆函数定理只保证小邻域可逆，不能给出整域的共同坐标。单连通性使无零函数可取对数与平方根；这些分支构造将产生足够多的有界单叶函数。

\subsection{单连通性}
\label{b3:subsec:050101}

本章单连通性沿用\subsecref{b3:subsec:020203}的闭路径收缩定义，区域必须非空、开且连通。凸区域、星形区域单连通；单位圆盘去掉一条从边界伸入内部的径向线段仍单连通，去掉一个内部点则不再单连通。

\ProseParagraphBreak

平面区域的拓扑判断不等于观察一幅图。一个常用的平面拓扑判据是：区域在扩充平面中的补集连通，当且仅当它单连通。该判据涉及 Jordan 分离与平面路径的多边形逼近，可作为本章所需的平面拓扑基础；实际分析证明仍以 Cauchy 定理的收缩条件为准。本章的构造不假定边界光滑，也不要求区域有界或边界只有一条简单曲线。

\ProseParagraphBreak

全平面本身单连通，却不能双全纯映到圆盘。否则圆盘坐标作为有界整函数由 Liouville 为常数，违反单射性。因此存在定理必须排除 \(U=\mathbb C\)，这一排除将具体用于在区域之外选择一点。

区域不必有简单边界。例如平面减去非负实轴是单连通的：极坐标可取 \(0<\arg z<2\pi\)，从而用 \(\log z\) 把它映到一条水平带，带是凸的。平面减去两个点却有两个独立绕行，绕其中任意一个点的闭路都不能收缩。对带有许多边界分支的区域，直接寻找一个能够连续选择的坐标，比只数图上的“孔”可靠。边界有尖点也不妨碍单连通性；是否能把共形映射连续延到尖点，属于第7章的另一问题。

\ProseParagraphBreak
单连通性只在区域内部讨论。开圆盘和闭圆盘不是同一对象，后者不是平面区域；开圆盘去掉一个边界点仍是原来的开圆盘。这个区分在谈到从边界伸入的裂缝时尤为必要。删去的线段应当确实含有内部点，其边界端只负责把障碍连接到区域外部，不能作为内部点参加闭路径。

\subsection{全纯对数}
\label{b3:subsec:050102}

\begin{theorem}[全纯对数]\label{b3:thm:05-holomorphic-log}
设 \(U\) 为单连通区域，\(f\) 在 \(U\) 上全纯且处处非零。则存在 \(U\) 上全纯函数 \(L\)，使 \(e^L=f\)。任意两个这样的函数相差 \(2\pi\mathrm{i}k\)，其中 \(k\in\mathbb Z\)。
\end{theorem}
\begin{proof}
函数 \(f'/f\) 全纯，由单连通 Cauchy 定理有原函数 \(H\)。求导得 \((f e^{-H})'=e^{-H}(f'-fH')=0\)，所以 \(fe^{-H}=c\ne0\)。选一个复数 \(\ell\) 满足 \(e^\ell=c\)，令 \(L=H+\ell\)，便有 \(e^L=f\)。若 \(L_1,L_2\) 都满足条件，差的指数为1，故差值落在离散集合 \(2\pi\mathrm{i}\mathbb Z\)。连续性与连通性使差为一个固定整数倍。
\end{proof}
给定 \(a\in U\) 和一个允许的 \(L(a)\)，构造中的常数选择就唯一。这个归一化对于沿路径继续分支很有用。

\ProseParagraphBreak

无零条件必不可少，因为指数函数从不为零。单连通性也不能直接删除：\(f(z)=z\) 在去心平面无零，却无全局对数。但它不是存在对数的必要条件，例如常值1在任何区域都有对数0；对某个特定函数的精确条件是 \(f'/f\) 的所有闭曲线周期为零。

\subsection{全纯平方根}
\label{b3:subsec:050103}

\begin{corollary}[全纯根]\label{b3:cor:05-holomorphic-root}
设 \(f\) 在单连通区域 \(U\) 上全纯且无零，\(m\ge1\) 为整数。则 \(f\) 有全纯 \(m\) 次根；任意两个根相差一个固定 \(m\) 次单位根因子。
\end{corollary}
\begin{proof}
取前一定理的 \(L\)，则 \(g=e^{L/m}\) 满足 \(g^m=f\)。两个无零根之商取值于有限集合 \(\{\zeta:\zeta^m=1\}\)，连续性与连通性使其恒定。
\end{proof}
若允许零点，局部 \(m\) 次根要求每个零点阶数能被 \(m\) 整除；例如 \(z\) 在含0的圆盘没有全纯平方根，\(z^2\) 却有。全局还需核对去掉零点后沿闭路的周期。下节只对无零函数取平方根，所以无需把这一更一般的根判别当作前提。

\ProseParagraphBreak

对 \(z-b\) 的平方根，若 \(b\notin U\)，得到 \(h^2=z-b\)。由 \(h(z_1)=h(z_2)\) 得 \(z_1=z_2\)，所以根本身单射。它的像与其负像不相交：若 \(h(z_1)=-h(z_2)\)，平方后点相同，再得 \(h(z_1)=0\)，与 \(b\notin U\) 矛盾。这个像集分离性质将产生有界函数。

\begin{proposition}\label{b3:prop:05-roots-with-zeros}
设 \(U\) 单连通，\(f\) 在 \(U\) 全纯且不恒零，\(m\ge2\)。则 \(f\) 有全纯 \(m\) 次根，当且仅当每个零点的阶数都能被 \(m\) 整除。
\end{proposition}
\begin{proof}
必要性由局部展开 \(f=g^m\) 立即得到。证明充分性时先去掉离散零点集 \(Z\)。对 \(U\setminus Z\) 中任意闭曲线，单连通 Cauchy 定理和有限奇点的留数公式给
\[
 \frac1{2\pi\mathrm{i}}\int_\gamma\frac{f'}f\,\mathrm{d}z
 =\sum_{a\in Z}\operatorname{Ind}(\gamma,a)\operatorname{ord}_a(f).
\]
右侧只有有限项且为 \(m\) 的整数倍。这里对于非简单曲线，可先按有限多边形分割成简单圈；小扰动避开零点后局部原函数保证积分不变。因此沿路径定义 \(g=c\exp(\frac1m\int f'/f)\)，其中 \(c^m=f(z_0)\)，闭路积分的指数均为1，得到单值全纯根。在零点 \(a\) 附近写 \(f=(z-a)^{mk}u\)，\(u\) 无零且在小盘有根 \(v\)。穿孔小盘上 \(g/((z-a)^k v)\) 是固定单位根，因此 \(g\) 全纯延过 \(a\)。
\end{proof}
与无零函数相比，证明多了两个步骤：先在删去零点的区域检查周期，再在每个零点局部补回。后一个步骤不能替代前一个步骤。在非单连通区域，所有零点阶数合适也可能没有根，例如去心平面的 \(z\) 没有零点，却无平方根。

\subsection{单叶函数}
\label{b3:subsec:050104}

区域上的单射全纯函数称为\term[danye-hanshu]{单叶函数}[univalent function]。单叶函数的导数不为零，由第一章局部标准形可知；其逆在像区域上全纯，因此是双全纯映射。本卷的“共形映射”若表示两个区域之间的等价，明确指这种双全纯映射，而局部共形函数不必全局单叶。

\ProseParagraphBreak

例如 \(e^z\) 导数处处非零，在宽度小于 \(2\pi\) 的水平带内单叶，却在整个平面不单叶。函数 \(z^2\) 在右半平面单叶，在包含0的圆盘不单叶，而且在0导数为零。判断单叶性需要比较任意两点，不能只核对导数。

\ProseParagraphBreak

若单叶函数 \(f:U\rightarrow\mathbb D\) 在某点 \(a\) 已确定值，可用 \(\phi_{f(a)}(w)=(w-f(a))/(1-\overline{f(a)}w)\) 将其送到0，再乘单位复数使导数成为正实数。这种归一化消除圆盘自同构中的自由度，不改变像是否为整个圆盘。
